% capitulos/cap5/seguranca.tex -- §5.12: segurança (ameaças, cabeçalhos e CSP, vulnerabilidades em dependências).
% Parte do Capítulo 5 (capitulos/cap5_arquitetura.tex); os rótulos \label e os números de seção não mudaram.

\section{Segurança}\label{sec:arq-seguranca}

A Tabela~\ref{tab:ameacas} relaciona as ameaças consideradas aos controles previstos, tomando como referência o OWASP Top 10 \cite{owasptop10} e as práticas de verificação do ASVS \cite{owaspasvs}.

\begin{longtable}{@{}L{3.3cm}L{4.4cm}L{5.4cm}@{}}
\caption{Ameaças e controles}\label{tab:ameacas}\\
\toprule
\textbf{Ameaça} & \textbf{Vetor} & \textbf{Controle}\\
\midrule
\endfirsthead
\toprule
\textbf{Ameaça} & \textbf{Vetor} & \textbf{Controle}\\
\midrule
\endhead
\bottomrule
\endfoot
XSS armazenado & Markdown com HTML ou URL \texttt{javascript:} & Sem HTML bruto (RN11); manter a transformação de URL do \arq{react-markdown}; revisar qualquer \en{plugin} novo; a CSP é uma segunda camada: barra \en{scripts} de outras origens e \texttt{eval}, mas não um \en{script} inline injetado, por causa do \texttt{'unsafe-inline'} (Seção~\ref{sec:arq-csp}, \adr{ADR-022})\\[1pt]
CSRF & Requisição de outro site com o \en{cookie} do administrador & \en{Cookie} \texttt{SameSite=Lax}; verificação de origem nas mutações (um \texttt{Origin} diferente do próprio site, ou \texttt{null}, responde \texttt{401} \arq{UNAUTHENTICATED} e nada é gravado; \adr{ADR-018}); rotas de escrita aceitam apenas JSON\\[1pt]
Injeção de SQL & Entradas em consultas & Prisma parametriza; proibido \texttt{\$queryRawUnsafe}\\[1pt]
Força bruta e enumeração & Tentativas repetidas de login & No máximo cinco tentativas por minuto por IP (resposta 429); mensagem genérica de erro, igual para e-mail inexistente e senha errada (UC08, CT02); o \en{proxy} deve definir \texttt{x-forwarded-for} (Seção~\ref{sec:arq-auth})\\[1pt]
Acesso indevido à API & Chamada direta sem passar pelo painel & Sessão verificada em cada Route Handler (RF21), antes de ler o corpo; Proxy apenas otimista; desde o M16, autorização por dono depois da sessão\\[1pt]
Rotas padrão do better-auth & Rotas de \arq{/api/auth/*} além do login (por exemplo, troca de senha e listagem de sessões) & Hoje exigem sessão e passam pela mesma limitação de requisições; o cadastro público está desabilitado (RF22). Desde o M16, todas fora de \texttt{get-session}, \texttt{sign-in/email} e \texttt{sign-out} respondem 404 (CT38; \adr{ADR-029})\\[1pt]
Acesso a recurso de outro dono & Trocar o identificador em uma rota de conteúdo, aparência, usuário ou notificação & \texttt{authorize} em toda rota, com a matriz falhando fechada; \texttt{authorId} da sessão; lote de reordenação recusado por inteiro; o banco impõe o dono único da árvore (CT35, CT37, CT39; RN19)\\[1pt]
Escalada de papel & Usuário que altera o próprio \texttt{role} ou \texttt{handle}, ou usa rotas do \en{plugin} \texttt{admin} & Campos do papel e do \en{handle} com \texttt{input: false}; sem \texttt{auth.api.*} sem \texttt{headers}; rotas \arq{/api/users} só para o administrador (CT37, CT38)\\[1pt]
Conta suspensa ou senha provisória & Sessão e \en{cookie} continuam valendo; senha provisória vazada ou nunca trocada & Sem \en{cache} de \en{cookie}; suspensão revoga as sessões e vence sozinha; senha exibida uma vez, nunca registrada, e troca obrigatória antes de qualquer outra ação (CT40, CT41; RN21, RN23)\\[1pt]
Enchimento do disco & Milhares de URLs inexistentes geram páginas em \en{cache} (cerca de 12 a 16~kB e 7 arquivos cada) & Proxy reescreve primeiros segmentos desconhecidos para uma página dinâmica que responde 404 (CT45); documento inexistente em espaço válido é risco aceito, com limite de taxa no \en{proxy} reverso e monitoramento do disco (Apêndice~\ref{ap:implantacao}; \adr{ADR-030})\\[1pt]
Abuso de escrita & Um usuário cria conteúdo ou corpos enormes & Cotas por usuário (500 conteúdos; 64~KiB em \emph{bytes}, ADR-033) e orçamento de renderização de 2~s na gravação, 413 antes de ler o corpo e limite de escritas por usuário (RN25; CT39)\\[1pt]
Abuso do administrador & Edição indevida da aparência ou leitura das notificações de outro usuário & Mínimo privilégio (RN19, aprovado em 04/10/2026); toda ação sobre outro usuário exige o motivo e notifica o dono (RN26; CT51)\\[1pt]
Exposição de segredos & \arq{.env} no repositório, na imagem Docker ou variáveis no navegador & \arq{.env} ignorado pelo Git e excluído da imagem por \arq{.dockerignore}, fornecido ao contêiner em tempo de execução; nenhuma variável com prefixo \arq{NEXT_PUBLIC_} para segredos\\[1pt]
Vazamento de rascunho & Nó \texttt{DRAFT} acessível por URL, pela árvore ou por página em cache & Filtro de publicação na camada de dados (RN15), que também oculta o espaço de usuário suspenso (RN23); revalidação do espaço ao alterar o estado (RN10, RN16); testes CT31, CT32 e CT43\\[1pt]
Link social malicioso & URL \texttt{javascript:} ou \texttt{data:} em um dos seis campos do rodapé & Validação de protocolo \texttt{http}/\texttt{https} no servidor (RN17); atributos \texttt{noopener} e \texttt{noreferrer}; teste CT33\\[1pt]
Injeção de CSS & Valor de cor malicioso & Validação \texttt{\#RRGGBB} no servidor (RN09)\\[1pt]
Dependências vulneráveis & Pacotes desatualizados, com vulnerabilidade conhecida ou maliciosos & Revisão de cada nova dependência pelo Arquiteto; \texttt{npm audit} a cada entrega; \en{lockfile} versionado; os alertas atuais foram avaliados e aceitos (Seção~\ref{sec:arq-npm-audit}, \adr{ADR-019} e \adr{ADR-032})\\[1pt]
\en{Clickjacking} & Página incorporada em \texttt{iframe} & CSP \texttt{frame-ancestors 'none'} e \arq{X-Frame-Options: DENY} (Seção~\ref{sec:arq-csp})\\[1pt]
Exclusão acidental & Cascata indevida & Confirmação explícita; \en{backup} (RNF10)\\[1pt]
Perda do banco & Remoção do diretório de dados, falha do disco da VPS ou atualização que descarta o volume & Banco em \en{bind mount} fora da imagem; \en{backup} periódico com \texttt{VACUUM INTO} e cópia externa (Apêndice~\ref{ap:implantacao}; risco R12)\\[1pt]
Arquivo de imagem hostil (épico v1, alvo) & SVG com \en{script}, poliglota (JPEG que também é HTML), EXIF com GPS, bomba de descompressão, arquivo gigante, nome de arquivo com caminho & Só PNG, JPEG e WebP; \emph{reencodificação} no servidor (sem metadados); limites de 2~MB e 25 megapixels; \texttt{sharp} fora do laço de eventos; o nome do arquivo nunca é usado; \texttt{/media} com tipo fixo e \texttt{nosniff} (Seção~\ref{sec:arq-upload}; \adr{ADR-034}; RN29)\\[1pt]
Abuso de CPU e de disco pelo envio (épico v1, alvo) & Envios em rajada ou simultâneos, trocas repetidas & Um processamento por vez por usuário e dois no servidor (503 \arq{BUSY}), limite de taxa próprio, troca que apaga a imagem anterior na mesma transação (sem imagem órfã) (Seção~\ref{sec:arq-upload})\\[1pt]
Troca de imagem alheia (épico v1, alvo) & Usuário envia ou remove a imagem de outro; usuário comum edita o rodapé ou a aparência do sistema & Autorização por dono nas rotas de imagem e de rodapé (403, RN19); o administrador só \emph{remove} o avatar de outro, com motivo e notificação (RN26)\\[1pt]
\end{longtable}
\fonte{Elaborado pelos autores.}

\subsection{Custo de Funções que Tocam Texto do Usuário}\label{sec:arq-custo-entrada}

O servidor é \emph{um só processo} Node.js, e o corpo de um conteúdo podia ter até 256~KiB (RN25; o padrão é de 64~KiB desde a \adr{ADR-033}): uma função que processa esse texto com custo superlinear deixa qualquer autor travar o servidor para todos. Por isso \textbf{o custo é um requisito de segurança}, e não de estilo. A regra, aplicada à \texttt{plainDescription} do SEO (merge M-029, R-018): (i) só se analisa o \emph{começo} do texto que a função precisa (\texttt{ANALYSIS\_LIMIT} de 4000 caracteres, para uma descrição de cerca de 160); (ii) nenhuma expressão regular tem quantificador ilimitado entre delimitadores que podem faltar: todos têm teto (\texttt{SPAN} de 300; recuo de no máximo 50), e o bloco de código cercado sai por uma varredura linear de linhas; (iii) o custo é \emph{medido com entradas hostis}, e não só com as felizes (\arq{tests/core/seo-cost.test.js}: 35 entradas de 256~KiB e janelas largas em que o sem-teto leva de 1 a 4~s e o limitado, de 20 a 150~ms; o ensaio \arq{doc/revisao/redos-seo.mjs}). \emph{Medido} (Node 24): antes da correção, 256~KiB de \texttt{[} repetido bloqueavam o processo por 128~s (e outras entradas, de 6 a 40~s); depois, os 20 casos do Arquiteto levam 12~ms ou menos.

\textbf{O renderizador de Markdown (P-016; decidido na \adr{ADR-033}; integrado no merge M-032, A-021).} O mesmo problema existe no \emph{pipeline} de renderização (\texttt{unified} com \texttt{remark-gfm}, \texttt{remark-math}, \texttt{rehype-katex} e \texttt{rehype-highlight}), que o servidor executa a cada geração de uma página pública e o navegador, a cada digitação no editor. Medido (parse e transformação, sem React; Node 24): um documento \emph{comum} de 256~KiB (títulos, listas, tabelas, código e fórmulas) leva cerca de 5,9~s, e um de 64~KiB, 0,8~s; entradas hostis de 64~KiB levam de 0,7 a 4,6~s (\texttt{> } repetido; \texttt{\$a } repetido, 2,4~s; \texttt{\$\$ }, 1,8~s; a crase, 1,1~s) e, de 256~KiB, até 30~s (\texttt{[a](} repetido) ou 13,6~s (\texttt{\$a }); o aninhamento profundo de citações estoura a pilha (\texttt{RangeError}) já com 64~KiB. A decisão do Engenheiro Chefe (06/10/2026, ADR-033) é baixar o padrão do limite do corpo para 64~KiB (RN25) \emph{e} recusar na gravação, com 422 \arq{CONTENT_TOO_COMPLEX}, o texto que não renderiza em 2~s, medido em um \emph{thread} separado (o aninhamento profundo, que estoura a pilha, cai na mesma recusa). \emph{Implementação} (\arq{src/lib/render-budget.js} e \arq{src/lib/render-worker.mjs}): o \emph{worker} executa o \emph{mesmo} \emph{pipeline} (a configuração é importada de \arq{markdown-pipeline.js}, nunca copiada); dura entre as chamadas (subir um custa cerca de 0,85~s, que não conta no orçamento do autor, e um por chamada somaria cerca de 1~s a toda gravação) e é encerrado e recriado quando o tempo acaba, quando morre ou depois de 5~min parado; tem teto de memória (512~MB de heap e 4~MB de pilha), de modo que um texto patológico derruba o \emph{worker}, e não o servidor; uma validação por vez, com três esperando (a seguinte recebe 503 \arq{BUSY} com \texttt{Retry-After: 3}); qualquer falha de \emph{renderização} (tempo, pilha, exceção do KaTeX, morte do \emph{worker}) é ``complexo demais'' (422), e a falha de \emph{subir} o \emph{worker}, que não é do autor, é 500. \emph{Medido} (Node 24, \emph{worker} quente, 62~KiB): documento comum, mediana de 106~ms e p95 de 142~ms; documento denso (fórmula, tabela e código em toda seção), 1,0~s e 1,3~s, isto é, 1,5$\times$ de folga sobre os 2~s; \texttt{GET /login} durante uma validação, de 18 a 30~ms; no ensaio HTTP, sete validações simultâneas deram 4 respostas 422 e 3 respostas 503. Com 64~KiB, só \texttt{> } repetido (5,6~s, que termina em \texttt{RangeError}) e \texttt{\$a } repetido (2,9~s) passam de 2~s na máquina de desenvolvimento; \texttt{\$\$ } (1,7~s), a crase (1,6~s) e \texttt{[a](} (0,9~s) são aceitos. \emph{Riscos residuais, aceitos e a reavaliar com a medição em produção:} (1) o orçamento de 2~s vale na gravação; um texto aceito pode custar até esse tempo, \emph{no processo principal}, a cada regeneração da página (ISR), e a geração não tem orçamento; (2) a fila é única, mas \emph{mitigada} (A-022, merge M-035): no máximo uma validação por usuário, em andamento ou na fila (a segunda do mesmo usuário recebe 503 \arq{BUSY}), liberada por qualquer término, inclusive a morte do \emph{worker}; um usuário hostil não impede mais os outros (medido no navegador: seis envios hostis da mesma conta deram 1 resposta 422 e 5 respostas 503, e a gravação normal de outro usuário foi atendida), mas quem fica atrás de uma validação hostil espera até cerca de 4~s, e quem tem várias contas (criadas só pelo administrador) ainda pode ocupar as quatro vagas da fila; o processador sai do \emph{worker} para \arq{src/lib/render-processor.mjs}, e testes provam que ele monta, junta por junta, o mesmo \emph{pipeline} do renderizador e que o \emph{worker} real recusa, com um orçamento curto, as entradas que só são caras \emph{com} o \emph{plugin} (matemática, KaTeX, realce), que as mutações do Arquiteto M60 e M61 tinham deixado escapar; (3) os corpos gravados antes do limite (até 256~KiB) continuam legíveis e são renderizados sem orçamento; como não há dados em produção antes da instalação, não há varredura; (4) o primeiro salvamento depois da partida paga a subida do \emph{worker}; (5) o \emph{navegador} também executa o \emph{pipeline}: o editor renderizava a pré-visualização mesmo oculta, e um texto hostil congelava a aba do \emph{próprio} autor (12~KB de \texttt{![} aninhado levavam 3,3~s, e 64~KiB, mais de 40~s), embora o servidor o recusasse em 2~s. \emph{Mitigado} (A-023, merge M-037): a pré-visualização oculta não renderiza nem checa nada; a visível espera 400~ms de pausa na digitação e só renderiza o texto que um \emph{Web Worker} aprovou em 1~s, executando o \emph{mesmo} processador do servidor (\arq{render-processor.mjs}); estourado o orçamento, o \emph{worker} é encerrado e a região diz ``A pré-visualização não está disponível: o texto é complexo demais''. Medido no Chromium sobre o \en{build} de produção: com 12~KB de \texttt{![} aninhado, o campo aceita a digitação em 38~ms e a mensagem aparece em 1,9~s; com 64~KiB de \texttt{*}, em 0,2~s; nenhuma tarefa longa; a pré-visualização oculta não cria \emph{worker}; a CSP não mudou (o \emph{worker} é um arquivo de \arq{/_next/static/}, do próprio site, coberto por \texttt{script-src 'self'}). O \emph{worker} do navegador usa um \emph{shim} mínimo (o Next resolve a condição ``browser'' dos pacotes também para o \emph{worker}, e duas \emph{builds} do \emph{pipeline} tocam o DOM ao carregar), de modo que ele deixa de decodificar entidades nomeadas e de converter o HTML do KaTeX em árvore, duas etapas lineares e baratas; o resto roda de verdade. (6) \emph{Risco residual novo, adiado para a v1.1:} o orçamento mede o \emph{processamento}, e não o custo de estilo e de DOM que o navegador paga para exibir o resultado; um documento \emph{legítimo} denso em fórmulas custa na \emph{thread} principal 0,5~s (8~KiB), 3,0~s (20~KiB) e 10,8~s (39~KiB), enquanto o \emph{worker} leva de 0,15 a 0,28~s; não é regressão nem texto hostil, mas o orçamento não o pega (um limite pelo peso do DOM gerado foi \emph{adiado para a v1.1} por decisão do Engenheiro Chefe em 08/10/2026; as opções são aceitar o risco, tentar antes um ajuste de CSS com \texttt{content-visibility: auto}, ou contar os elementos que o processador gera e recusar, com o mesmo 422 \arq{CONTENT_TOO_COMPLEX}, o texto acima de um teto calibrado por medição, o mesmo \arq{render-processor.mjs} no servidor e no navegador).

\subsection{Envio de Imagens e Cookies Novos (Épico v1)}\label{sec:arq-upload}

\textbf{Situação: integrado (A-025 especificou; A-029, merge M-042, 09/10/2026; o Arquiteto o reexecutou com as mutações M97 a M132 e o ensaio HTTP contra o \en{build} de produção).} O envio de imagem é a única superfície de ataque nova do épico (\adr{ADR-034}). A defesa é \emph{transformar}, e não só validar: a imagem que o servidor guarda e serve é a que o \texttt{sharp} gerou, e não a que o usuário enviou. Controles, na ordem em que atuam:
\begin{enumerate}
  \item \textbf{Autenticação, origem e papel} (Seção~\ref{sec:arq-autorizacao}); a conta com a senha provisória pendente não envia.
  \item \textbf{Leitura com limite}: o corpo \texttt{multipart} é lido até 2~MB mais a folga fixa de 16~KiB das fronteiras e dos cabeçalhos; acima, 413, \emph{antes} de decodificar, e o arquivo em si também é conferido (2~MB${}+{}$1 byte é 413); \texttt{Content-Type} diferente de \texttt{multipart/form-data}, ou um corpo malformado, é 400.
  \item \textbf{Concorrência e taxa}: um processamento por vez por usuário e dois no servidor (503 \arq{BUSY}); limite de taxa próprio (429); o \texttt{sharp} roda fora do laço de eventos do servidor.
  \item \textbf{Decodificação e reencodificação}: aceita PNG, JPEG e WebP (a decisão vem do \emph{conteúdo}, isto é, da assinatura dos primeiros bytes; o nome do arquivo nunca é usado, e o tipo declarado só é conferido: um tipo declarado diferente do conteúdo é recusado, e o vazio ou \texttt{application/octet-stream} conta como não declarado); a dimensão é lida do cabeçalho, sem decodificar, e acima de 25 megapixels (25.000.000 de pixels; o limite exato é aceito) é recusada, com o \texttt{limitInputPixels} do \texttt{sharp} como segunda trava; saída no formato e no tamanho do alvo (avatar 256$\times$256 WebP; logotipo até 512~px, WebP; \en{favicon} 64$\times$64, PNG), \emph{sem metadados} (o EXIF com GPS não sai); SVG, GIF, HTML ou ZIP renomeado, poliglota, arquivo vazio e arquivo corrompido terminam em 422 \arq{INVALID_IMAGE}, sem eco do arquivo.
  \item \textbf{Gravação}: a imagem nova, o ponteiro do dono e a exclusão da antiga acontecem numa só transação; o identificador é novo, de modo que a URL imutável nunca mostra a imagem errada.
  \item \textbf{Servir}: \texttt{GET /media/<id>} é público; \texttt{Content-Type} fixo da saída, \texttt{nosniff}, \texttt{ETag}, \texttt{Cache-Control} imutável; 404 sem distinguir. A imagem de um espaço suspenso continua servida (\adr{ADR-034}, item 8): o identificador é opaco e só as páginas do espaço, que dão 404, o referenciam.
\end{enumerate}
A CSP \emph{mantém} \texttt{img-src 'self' data: https:} na v1 (Tabela de diretivas, adiante; \adr{ADR-034}, item 11): logotipo, \en{favicon} e avatar são sempre servidos por \texttt{/media}, da mesma origem, mas as imagens externas dos documentos Markdown continuam carregando. O \texttt{sharp} é uma dependência nativa: entra na avaliação de \texttt{npm audit} a cada entrega (Seção~\ref{sec:arq-npm-audit}), e o CT29 é repetido com ele (Apêndice~\ref{ap:implantacao}).

\textbf{Cookies novos.} \texttt{vitrine-modo} (\texttt{light} ou \texttt{dark}; RN32) e \texttt{vitrine-logado} (valor \texttt{1}; RN34) são \emph{funcionais}, não carregam dado pessoal nem segredo e não têm \texttt{HttpOnly}, porque o navegador precisa escrevê-los e lê-los: o primeiro só muda o aspecto da página; o segundo é só uma dica para a interface decidir se consulta a sessão e \emph{nunca autoriza nada} (a autorização continua no servidor, em cada rota e página protegida). Ambos usam \texttt{Path=/}, \texttt{SameSite=Lax} e \texttt{Secure} sob HTTPS; o do modo dura um ano, e o marcador é apagado ao sair e quando a consulta de sessão devolve vazio. Um marcador forjado só faz o navegador de quem o forjou fazer uma consulta de sessão a mais.

\subsection{Cabeçalhos de Segurança e Política de Segurança de Conteúdo}\label{sec:arq-csp}

Além dos controles acima, o Vitrine envia cabeçalhos de segurança em \emph{todas} as respostas --- páginas, API e arquivos estáticos. Eles estão definidos em \arq{security-headers.mjs} (na raiz do repositório; módulo puro, sem imports, testado sem o Next.js) e são aplicados pela função \texttt{headers} do \arq{next.config.mjs}, que também desliga o \texttt{X-Powered-By}. A decisão e as alternativas estão na \adr{ADR-022}. A Tabela~\ref{tab:cabecalhos} lista os cabeçalhos, e a Tabela~\ref{tab:csp-diretivas}, as diretivas da \en{Content Security Policy} (CSP).

\begingroup\small
\begin{longtable}{@{}L{3.6cm}L{4.4cm}L{5.2cm}@{}}
\caption{Cabeçalhos de segurança enviados}\label{tab:cabecalhos}\\
\toprule
\textbf{Cabeçalho} & \textbf{Valor} & \textbf{Finalidade}\\
\midrule
\endfirsthead
\toprule
\textbf{Cabeçalho} & \textbf{Valor} & \textbf{Finalidade}\\
\midrule
\endhead
\bottomrule
\endfoot
\arq{X-Content-Type-Options} & \texttt{nosniff} & O navegador não interpreta uma resposta com tipo diferente do declarado\\
\addlinespace
\arq{Referrer-Policy} & \texttt{strict-origin-when-\allowbreak cross-origin} & Outros sites recebem só a origem, e não a URL completa\\
\addlinespace
\arq{X-Frame-Options} & \texttt{DENY} & Impede o enquadramento (\en{clickjacking}) em navegadores que ignoram \texttt{frame-ancestors}\\
\addlinespace
\arq{Permissions-Policy} & \texttt{camera=(), microphone=(), geolocation=(), payment=(), usb=()} & Desliga recursos do navegador que o Vitrine não usa\\
\addlinespace
\arq{Content-Security-Policy} & Tabela~\ref{tab:csp-diretivas} & Restringe a origem de cada tipo de recurso; enviado \emph{somente em produção}\\
\addlinespace
\emph{Não enviados} & \arq{Strict-Transport-Security} e \texttt{upgrade-insecure-requests} & TLS, HSTS e o redirecionamento para HTTPS são do \en{proxy} reverso (Apêndice~\ref{ap:implantacao})\\
\end{longtable}
\fonte{Elaborado pelos autores, a partir de \arq{security-headers.mjs} (merge M-008).}
\endgroup

\begingroup\small
\begin{longtable}{@{}L{2.8cm}L{3.3cm}L{7.1cm}@{}}
\caption{Diretivas da CSP adotada e justificativas (ADR-022)}\label{tab:csp-diretivas}\\
\toprule
\textbf{Diretiva} & \textbf{Valor} & \textbf{Justificativa}\\
\midrule
\endfirsthead
\toprule
\textbf{Diretiva} & \textbf{Valor} & \textbf{Justificativa}\\
\midrule
\endhead
\bottomrule
\endfoot
\texttt{default-src} & \texttt{'self'} & Ponto de partida restritivo: o que não tem diretiva própria só carrega do próprio site\\
\addlinespace
\texttt{script-src} & \texttt{'self' 'unsafe-inline'} & \textbf{Primeiro uso de \texttt{'unsafe-inline'}.} O Next.js emite \en{scripts} inline de hidratação em todas as páginas; sem ele, a página não hidrata e o formulário de login cai no envio nativo do navegador, por GET, com o e-mail e a senha na URL. Sem \texttt{'unsafe-eval'} e sem \en{nonce}\\
\addlinespace
\texttt{style-src} & \texttt{'self'} & Folhas de estilo só do próprio site (os arquivos do \en{build}); folhas e elementos \texttt{<style>} inline continuam proibidos\\
\addlinespace
\texttt{style-src-attr} & \texttt{'unsafe-inline'} & \textbf{Segundo uso.} Só para atributos \texttt{style}: o tema do banco vai no atributo \texttt{style} do \texttt{<html>} (Seção~\ref{sec:arq-tema-estatico}), e as fórmulas do KaTeX também os usam; atributos de estilo não aceitam \en{nonce}. Sem a diretiva, a paleta do banco é ignorada em silêncio\\
\addlinespace
\texttt{img-src} & \texttt{'self' data: https:} (vigente e mantido na v1) & Vigente: as imagens do Markdown são URLs externas (D-10); \texttt{http:} é bloqueado. \emph{Épico v1 (\adr{ADR-034}, item 11, decisão do Engenheiro Chefe em 08/10/2026):} o \texttt{https:} \textbf{permanece}; o \texttt{'self'} cobre \texttt{/media/<id>} (logotipo, \en{favicon} e avatar, nunca por URL externa), e o \texttt{data:} é dos ícones inline e da pré-visualização local do envio. Nenhum componente \emph{novo} carrega imagem externa; o \texttt{MarkdownRenderer} não muda. \emph{Plano da v1.1:} upload de imagens de conteúdo e CSP sem \texttt{https:}\\
\addlinespace
\texttt{font-src} & \texttt{'self'} & As fontes do site e do KaTeX devem vir empacotadas no \en{build}: a folha do KaTeX é importada pelo \en{bundle}, e não de uma CDN\\
\addlinespace
\texttt{connect-src} & \texttt{'self'} & A interface só chama a própria API\\
\addlinespace
\texttt{object-src} & \texttt{'none'} & Sem plugins (\texttt{object} e \texttt{embed})\\
\addlinespace
\texttt{base-uri} & \texttt{'self'} & Impede trocar o \texttt{<base>} para redirecionar URLs relativas\\
\addlinespace
\texttt{form-action} & \texttt{'self'} & Formulários só enviam dados ao próprio site\\
\addlinespace
\texttt{frame-ancestors} & \texttt{'none'} & Nenhuma página pode ser incorporada em \texttt{iframe} (\en{clickjacking})\\
\end{longtable}
\fonte{Elaborado pelos autores, a partir de \arq{security-headers.mjs} e do guia de CSP do Next.js 16 \cite{nextjsdocs}.}
\endgroup

\textbf{Por que não há \en{nonce}.} Segundo o guia de CSP do Next.js 16 \cite{nextjsdocs}, o \en{nonce} é gerado a cada visita e só pode ser aplicado na renderização no servidor: todas as páginas teriam de ser dinâmicas, a otimização estática e a ISR (Seção~\ref{sec:arq-isr}) ficariam desligadas e as páginas públicas deixariam de ser servidas do cache, o que contraria a \adr{ADR-005} e o RNF02. O Vitrine adota, então, a forma que o próprio guia indica para quem mantém a geração estática, porém mais estrita que o exemplo dele: as folhas de estilo continuam restritas a \texttt{'self'}, e apenas os atributos \texttt{style} são liberados, por \texttt{style-src-attr}. Os \emph{hashes} com Subresource Integrity, que dispensariam o \texttt{'unsafe-inline'} dos \en{scripts}, ainda são experimentais no Next.js 16.

\textbf{O que a CSP protege e o risco aceito.} Com \texttt{script-src 'unsafe-inline'}, a CSP \emph{não} detém um \en{script} inline injetado; a defesa primária contra XSS continua sendo a ausência de HTML bruto (RN11), a transformação de URLs do \arq{react-markdown} (que recusa \texttt{javascript:}), o escape do React e a validação dos links sociais (RN17). A CSP é a segunda camada: barra \en{scripts} de outras origens, \texttt{eval}, folhas de estilo inline, imagens em \texttt{http:}, conexões e formulários para outros sites, plugins, troca do \texttt{<base>} e enquadramento em \texttt{iframe}. O risco residual de \texttt{img-src https:} é uma imagem de rastreamento de qualquer domínio, aceito porque o conteúdo referencia imagens externas (D-10); a v1 não o altera (\adr{ADR-034}, item 11), e a v1.1 o remove com o upload de imagens de conteúdo.

\textbf{Somente em produção e como verificar.} O \texttt{next dev} precisa de \texttt{'unsafe-eval'} (o React o usa para depuração), que a política não concede; por isso a CSP \emph{não} é enviada fora de produção, e um recurso novo pode funcionar no desenvolvimento e ser bloqueado em produção. Todo recurso novo que carregue \en{script}, estilo, fonte, imagem, quadro ou conexão (o editor e o KaTeX, as páginas públicas, o rodapé) deve ser verificado no navegador sobre \texttt{npm run build} e \texttt{npm run start}, com o console aberto: o formulário de login hidrata, a paleta do banco é aplicada e não aparece nenhuma mensagem de violação de CSP. A verificação do merge M-008 está na Seção~\ref{sec:testes-seguranca}.

\textbf{O Zod no navegador e o \texttt{eval}.} Achado do merge M-011, só visível no \en{build} de produção: ao criar o primeiro esquema de objeto, o Zod 4 sonda \texttt{new Function("")} para saber se pode compilar um validador mais rápido (JIT). O Zod captura o erro e segue, mas a CSP, sem \texttt{'unsafe-eval'}, faz o navegador registrar a tentativa como violação (o evento \texttt{securitypolicyviolation} e um erro no console), a cada carga. O editor chama \texttt{z.config(\{ jitless: true \})} em \arq{zod-config.js}, que o \arq{schema.js} importa \emph{antes} do \arq{validators.js} (a ordem de avaliação dos módulos é o que garante a configuração a tempo); a validação dá o mesmo resultado, sem o atalho compilado. A configuração vale só para o código do navegador: o servidor não é afetado. Quem importar o \arq{validators.js} em outro código de navegador deve configurar o Zod antes dele. Um teste (\arq{tests/ui/editor-zod-csp.test.js}) registra toda construção de função a partir de texto e prova que a sonda não acontece.

\textbf{Alterações.} Uma nova origem ou um novo tipo de recurso é um \texttt{[PEDIDO]} ao Core, exige a verificação acima e muda o teste que confere as onze diretivas uma a uma (\arq{tests/core/infra.test.js}); enfraquecer a política exige uma nova ADR. Vídeos e quadros incorporados, fontes e \en{scripts} de CDN e ferramentas de análise de acesso ficam fora do escopo desta versão. HSTS e TLS ficam a cargo do \en{proxy} reverso, que também não deve substituir nem remover os cabeçalhos da aplicação (Seção~\ref{sec:proxy}).

\subsection{Vulnerabilidades Conhecidas em Dependências}\label{sec:arq-npm-audit}

Em 01/10/2026, na revisão do merge M-004, o \texttt{npm audit} acusou \textbf{quatro vulnerabilidades de severidade alta}, todas em dependências transitivas do Prisma 7 (Tabela~\ref{tab:npm-audit}). O npm só as corrige com \texttt{npm audit fix -{}-force}, que instalaria o \texttt{prisma@6.19.3}: um rebaixamento de versão maior, incompatível com o que o projeto adota (\arq{prisma.config.ts}, gerador \texttt{prisma-client} e \en{driver adapter}; \adr{ADR-014}). A \texttt{7.10.0}, instalada, é a última versão 7 publicada, e o \texttt{npm audit} lista como afetada toda a faixa do \arq{prisma} até a \texttt{8.1.0-dev.6}; portanto, não há versão 7 corrigida. Decidiu-se \textbf{manter o Prisma 7 e não aplicar a correção automática}, aceitando os alertas, pelas razões da tabela (\adr{ADR-019}). Como \arq{prisma} agora está em \texttt{dependencies} (para aplicar migrações na imagem), esses pacotes também estão presentes na imagem de produção, embora não sejam exercidos.

\begingroup\small
\begin{longtable}{@{}L{3.3cm}L{5.4cm}L{5.5cm}@{}}
\caption{Alertas do \texttt{npm audit} aceitos em 01/10/2026}\label{tab:npm-audit}\\
\toprule
\textbf{Pacote} & \textbf{Aviso} & \textbf{Avaliação para o Vitrine}\\
\midrule
\endfirsthead
\toprule
\textbf{Pacote} & \textbf{Aviso} & \textbf{Avaliação para o Vitrine}\\
\midrule
\endhead
\bottomrule
\endfoot
\arq{deepmerge-ts} (antes da 8.0.0; usado por \arq{@prisma/config}) & Esgotamento de pilha ao mesclar grafos de objetos recursivos (GHSA-ggr8-5vv4-36mx) & Roda na CLI do Prisma, ao mesclar a configuração de \arq{prisma.config.ts}, escrita pela equipe; não recebe entrada de usuário nem é exercido pelo servidor\\
\addlinespace
\arq{@prisma/config} e \arq{prisma} & Alertas em cadeia, por dependerem do \arq{deepmerge-ts} e do \arq{mysql2} & O mesmo das linhas vizinhas; sem correção em nenhuma versão 7 disponível\\
\addlinespace
\arq{mysql2} (até a 3.23.0; dependência do \arq{prisma} e do \arq{better-auth}) & Rebaixamento do plugin de autenticação para \arq{mysql_clear_password}, com vazamento de credenciais (GHSA-3f6p-5ww8-9rcr), e descompressão (\en{inflate}) sem limite no protocolo comprimido, com negação de serviço (GHSA-rgwj-5xj2-c3m3) & Só é exercido ao conectar a um servidor MySQL; o Vitrine usa SQLite por \arq{better-sqlite3} e nunca abre conexão MySQL\\
\end{longtable}
\fonte{Elaborado pelos autores, a partir da saída do \texttt{npm audit} e de \texttt{npm ls} (merge M-004).}
\endgroup

A decisão será reavaliada a cada entrega (o \texttt{npm audit} é critério de saída, Seção~\ref{sec:testes-saida}) e sempre que o Prisma publicar uma versão 7 com as dependências corrigidas. Enquanto isso, os alertas aceitos não contam como vulnerabilidade crítica pendente.

\textbf{Complemento de 03/10/2026.} Na primeira execução do projeto em Linux (ambiente de nuvem do Arquiteto), o \texttt{npm audit} acusou \textbf{nove} alertas de severidade alta, e não quatro. Os quatro da Tabela~\ref{tab:npm-audit} continuam os mesmos; os cinco novos vêm de um único aviso, publicado depois, sobre o \arq{braces} (até a 3.0.3; esgotamento de pilha com padrões aninhados em excesso, GHSA-vfj7-8cjw-p6xm), que chega pela cadeia do \arq{eslint-config-next}, do \arq{@next/eslint-plugin-next}, do \arq{fast-glob}, do \arq{micromatch} e do \arq{braces}. É uma cadeia exclusiva de ferramenta de desenvolvimento (\texttt{devDependencies}): o servidor não a executa, e o \texttt{eslint} só avalia padrões escritos pela equipe. O npm só a corrige com \texttt{-{}-force}, que rebaixaria o \arq{eslint-config-next} para a 14.2.35, e não há versão corrigida do \arq{braces} na faixa que o \arq{micromatch} aceita. A comparação com a lista acima, que a \adr{ADR-019} exige justamente para que os alertas aceitos não mascarem os novos, mostrou o mesmo conjunto de nove alertas na \texttt{main}, no merge M-010 e no M-011: nenhum merge acrescentou alerta. O Engenheiro Chefe aceitou os cinco alertas em 04/10/2026, pelo mesmo critério da \adr{ADR-019} (pacote não exercido pelo produto), com uma condição verificável: o \arq{braces} só pode aparecer na cadeia de desenvolvimento. O \texttt{npm ls braces -{}-omit=dev} imprime a árvore vazia, e o \texttt{npm audit -{}-omit=dev} acusa só os quatro alertas do Prisma; a aceitação está na \adr{ADR-032}, e a reavaliação ocorre quando o \arq{eslint-config-next} ou o \arq{micromatch} publicarem versão corrigida, ou se a condição deixar de valer.
